% !PRINTED 178
\weberpage{178}
%<<CONT
soluten auch relative Invarianten existiren, und diese können zur
Vereinfachung des Formenproblems benutzt werden.

Ist $r$ eine solche relative Invariante, so wird eine gewisse
Potenz von $r$, deren Grad $e$ ein Theiler des Grades $\mu$ von $S$
ist, eine absolute Invariante, und $r$ wird also durch eine \gesperrt{reine
Gleichung} in $\Omega$ bestimmt.

Alle Substitutionen von $S$, durch die $r$ ungeändert bleibt,
bilden für sich eine Gruppe $T_r$, und die Gruppe $T_r$ ist, wie wir
in \S.\,40 gesehen haben, ein Normaltheiler von $S$.

Ferner aber sehen wir, dass jede Function der Variablen $x$,
die durch die Substitutionen der Gruppe $T_r$ ungeändert bleibt,
rational durch $r$ ausgedrückt werden kann, d.\,h. in dem durch
Adjunction von $r$ aus $\Omega$ abgeleiteten Körper $\Omega_r$ enthalten ist.

Nach \S.\,40 nämlich giebt es eine Substitution $E$ in $S$, und
eine primitive $e^{\text{te}}$ Einheitswurzel $\varepsilon$, so dass $r$ durch $E$ in $\varepsilon r$
übergeht, und alle Werthe, die $r$ annehmen kann:
\[
r, \varepsilon r, \varepsilon^2 r, \ldots \varepsilon^{e-1} r,
\]
erhält man durch Wiederholung der Substitution $E$. Bedeutet
ferner $\tau$ eine Function der $x$, die die Substitutionen der Gruppe
$T_r$ gestattet, und durch $E$ und seine Potenzen in
\[
\tau, \tau_1, \tau_2, \ldots \tau_{e-1}
\]
übergeht, so gestatten alle diese Functionen gleichfalls die
Substitutionen von $T_r$, weil $T_r$ ein Normaltheiler von $S$ ist. Die
Function
\weq{1}{(\varepsilon^{-h}, \tau) = \tau + \varepsilon^{-h}\tau_1 + \cdots + \varepsilon^{-h(e-1)}\tau_{e-1}, \qquad (h = 0, 1, 2, \ldots e-1),}
die nach Analogie der \gesperrt{Lagrange}'schen Resolventen (Bd.\,I,
\S.\,164) gebildet ist, nimmt durch Anwendung der Substitution $E$
den Factor $\varepsilon^h$ an, und wenn wir also
\weq{2}{(\varepsilon^{-h}, \tau) = r^h \Psi_h
}
setzen, so ist $\Psi_h$ eine absolute Invariante, also im Körper $\Omega$
enthalten.

Aus (1) und (2) ergiebt sich aber
\weq{3}{e\tau = \Psi + r\Psi_1 + r^2\Psi_2 + \cdots + r^{e-1}\Psi_{e-1},}
wodurch die Behauptung erwiesen ist.

Setzen wir
\weq{4}{\mu = e\nu,}
%CONT>>
